% !TEX root = ../main.tex
% capitulos/1_introducao.tex

\section{APRESENTAÇÃO OU INTRODUÇÃO}

Estas normas têm como objetivo dar uma orientação geral aos autores dos artigos no momento em que forem redigir e, principalmente, quando forem organizar e digitar seus artigos científicos.

Esse documento já está configurado com as normas pré-estabelecidas pelos editores da Revista Somma e, para segui-las, basta substituir os textos de descrição pelo conteúdo do artigo. O texto do artigo deve ser estruturado preferencialmente contemplando os seguintes itens: introdução, método, resultados e conclusão. Acrônimos e abreviações devem estar entre parênteses e serem precedidos de seu significado completo quando do primeiro uso no texto. 

Palavras estrangeiras devem ser grafadas em \textit{itálico}. Para ênfase ou destaque usar \textbf{negrito}, `aspas simples' ou ``aspas duplas''. Exemplo de paráfrase cuja autoria da fonte é parte do texto: De acordo com \textcite{santana2009} manuscritos submetidos devem se apresentar exatamente como solicitado no template da revista.

Exemplo de paráfrase cuja autoria da fonte não é parte do texto: manuscritos submetidos devem se apresentar exatamente como solicitado no template da revista \cite{santana2009}.

No caso de citações diretas curtas (até 3 linhas), as mesmas devem ser colocadas entre aspas duplas ``manuscritos submetidos aguardam em torno de três meses devido o processo editorial ser lento'' \cite[p. 14]{santana2009}.

No caso de citações longas (mais de três linhas), este é o exemplo. 
\begin{citacao}
Lorem ipsum dolor sit amet, consectetur adipiscing elit. Ut vulputate tincidunt turpis at tincidunt. degerenuim vitallis quereris suprimum totis la menera cual qieu vestapenin paenas eojlpsm de quiel fiquenopos acerbalium neolipis vulguoris dermates enxaques delamore qual tu seia esupendistance.
\end{citacao}

Recomenda-se nunca terminar uma seção com citação longa. Procure continuar com o texto de forma a estabelecer uma ligação com o item/seção seguinte.

\subsection{Normas para submissão de artigos}

Recomenda-se que o artigo tenha no mínimo 6 e no máximo 20 páginas na versão final publicada, contando as referências, figuras e tabelas. 

\subsubsection{Exemplo}
Segue abaixo um exemplo de organização do artigo em forma de tópicos, bem como a formatação de cada um.
